\documentclass[11pt]{article}
\usepackage[margin=1in]{geometry}
\usepackage{amsmath,amssymb,amsthm,mathtools,bm}
\usepackage{microtype}
\usepackage[T1]{fontenc}
\usepackage{lmodern}
\usepackage{booktabs,enumitem}
\usepackage[hidelinks]{hyperref}
\setlist{nosep}
\newtheorem{theorem}{Theorem}[section]
\newtheorem{proposition}[theorem]{Proposition}
\newtheorem{lemma}[theorem]{Lemma}
\theoremstyle{definition}
\newtheorem{assumption}[theorem]{Assumption}
\theoremstyle{remark}
\newtheorem{remark}[theorem]{Remark}
\newtheorem{conjecture}{Conjecture}

\newcommand{\Qbar}{\overline{Q}}

\title{\textbf{Synopsis: Early-Universe Specialization of Hysterical Instability}\\[0.35em]
\large Instability inheritance, rapid-formation reading, and the FLRW timing suite}
\author{Andrew Bruce Boyd}
\date{September 2026}

\begin{document}
\maketitle
\noindent\textit{Computational note.}
Proofs, exploratory experiments, and paper writing were assisted by generative AI.
Mathematical theorems stated in this paper are formally verified in Lean~4 in the
companion specifications, as faithful ports of the TeX arguments at the abstractions
recorded there, except results explicitly tagged as \emph{literature hypotheses}
(standard theorems cited from the literature and not re-proved here) or
\emph{physical inputs} (modeling assumptions outside mathematics).

\begin{abstract}
Short guide to the early-universe specialization note.
Instability and threshold times are inherited from the response-theory paper.
What is new is the cosmological closure, a dust-rate comparison, an interpretive reading as rapid early galaxy formation, and the full FLRW timing suite formerly developed as a follow-on numerical paper.
\end{abstract}

\section{Why this note exists}
The response-theory paper already proves closed-loop instability when $\alpha(J)>0$, and finite logarithmic threshold times for nonzero unstable projections.
This note checks that a minimal post-decoupling cosmological model lands in that supercritical phase, then asks whether the same loop can act fast enough to matter for the earliest galaxies.

\section{Cosmological closure}
Local reservoir linearization $\dot q+\gamma q=\sigma\delta$ with attractive dust closure produces a frozen Jacobian with
\[
P(\lambda)=(\lambda+\gamma)(\lambda^2+2H\lambda-A)-A\beta\sigma.
\]
Positive coupling gives $P(0)<0$, hence $\alpha(J)>0$, so the inherited instability package applies.
A response-coupled root also satisfies $\lambda_+>\lambda_b$ for ordinary dust.

\section{Interpretation}
Candidate mechanism for rapid early galaxy formation: baryonic seeds amplify a response sector that feeds back into gravity, without a CDM perturbative scaffold.
JWST has spectroscopically confirmed luminous galaxies at $z\simeq14$.

\section{FLRW timing suite}
Planck-like background; $x=\beta q$; $\kappa=\beta\sigma/H$, $g=\gamma/H$ held constant; baryon-only clustering; fiducial start at $z_{\mathrm{on}}=1089$ with $\delta_i=10^{-5}$.

\begin{proposition}[Timing viability at $z\simeq14$]
For $g=0$, $\kappa_{\mathrm{crit}}(z_{\mathrm{nl}}=14)\simeq129.8$.
\end{proposition}

Further numerical results retained from the timing suite:
\begin{itemize}
\item Loss raises $\kappa_{\mathrm{crit}}$ ($\simeq142$, $171$, $254$ for $g=0.3$, $1$, $3$).
\item Delayed activation raises it more strongly ($\simeq197$--$540$ for $z_{\mathrm{on}}=600$--$200$).
\item Near threshold, $x=\beta q$ amplifies ahead of $\delta$.
\item Smaller seeds cross later at fixed $\kappa$ (logarithmic timing, not weaker exponent).
\end{itemize}
This is timing possibility, not an abundance fit.
CMB transfer, stellar masses, and luminosity functions remain open.

\end{document}
